%
%============= CHAPTER BREAK =================%
%
\chapter{Philumenos (ca 180--220 CE)}
\vspace*{3mm}%
\justifying%
%
\noindent{}Philumenos the Greek medical author lived and wrote around the late 2nd to early 3rd centuries CE.\endnote{%
Bouras \& Zipser 2019:42, 96; Diller 1941:209 notes that Philumenos cites Archigenes, Soranus, and Marcellus, %
and makes use of the Pneumatist Herodotus, but is not cited by Galen.} \\[2mm] %

\noindent{}\textbf{ἐξέραμα} / \textit{ekserama} / “vomit” appears only one time in the NT (2Pt 2:22), %
and only 51 times in the TLG. %
Its first known appearance in any non-Christian text is in the treatise %
\textit{On Poisonous Animals and Their Remedies} / \textit{De venenatis animalibus eorumque remediis}. %
In the section on rabies, or what Greek medical authors called \textit{hydrophobia}, %
Philumenos records the following treatments: %

\vspace*{3mm} 
 \begin{quote}πρὸϲ ὑδροφόβαϲ τὸ διὰ πυτίαϲ. πυτίαϲ λαγωοῦ, Λημνίαϲ ϲφραγῖδοϲ, 
 ἀρκευθίδων, γεντιανῆϲ ἀνὰ {\greeksymbol 𐅻} δ, δαφνίδων, %
 ϲμύρνηϲ ἀνὰ {\greeksymbol 𐅻} η ἀναλάμβανε μέλιτι. %
 (13) ἡ δόϲιϲ κυάμου Αἰγυπτίου τὸ μέγεθοϲ. %
  ἄλλο· | \textbf{ἐξέραμα} κυνὸϲ ἄραϲ ξήρανον καὶ λεάναϲ θὲϲ ἐν πυξίδι %
  καὶ χρῶ ὡϲ καλλίϲτῳ πρὸϲ λυϲϲοδήκτουϲ θαρρῶν. (\textit{venen.anim.} 4.13)\\[1mm] %
 For rabies, the [remedy is made] of rennet: %
 rennet of a hare, with the seal of Lemnia, %
 juniper-berry, gentian root 4 drachms, bayberry, %
 myrrh 7 drachms, mix up with honey. %
 (13) The dose is the size of an Egyptian bean. %
 Another [remedy]: take the dried \textbf{vomit} of a dog and grind it, %
 put it in box-wood, and use as is best for those bitten by a rabid dog, %
 taking courage. %
 \end{quote} % 
 \vspace*{3mm}

\noindent{}The passage in 2 Peter shares several similarities, %
not just the new/rare lemma ἐξέραμα / \textit{ekserama} / “vomit”, %
but also its association with dogs, a nearby highly positive/optimistic reference %
(cp. “for the best” / καλ\-λίϲτῳ and “better” / κρεῖττον), %
and especially the vivid and disgusting image of eating vomit %
used in combination with the motif of unhealthy human contact with unclean/diseased dogs. %

\vspace*{1mm} 
 \begin{quote}κρεῖττον γὰρ ἦν αὐτοῖς
 μὴ ἐπεγνωκέναι τὴν ὁδὸν τῆς δικαιοσύνης ἢ ἐπιγνοῦσιν
 ὑποστρέψαι ἐκ τῆς παραδοθείσης αὐτοῖς ἁγίας ἐντολῆς.
 ²²  συμβέβηκεν αὐτοῖς τὸ τῆς ἀληθοῦς παροιμίας,
  Κύων ἐπιστρέψας ἐπὶ τὸ ἴδιον \textbf{ἐξέραμα},
 καί, Ὗς λουσαμένη εἰς κυλισμὸν βορβόρου. \\[1mm] %
 For it were better for them not to have recognized the way of the righteous than, %
 after recognizing it, to turn back from the holy command given over to them. %
 ²² The proverb of truth has come together for them: % 
 “A dog turning back to its own \textbf{vomit}”, %
 and, “A sow has been washed to roll in mud”. %
 \end{quote} % 
 \vspace*{3mm}

As Neyrey and other commentators have noted, %
in this slightly larger segment (2:20--22), 2 Peter departs from its source in Jude.\endnote{%
Neyrey 2021:220.} %
Antecedents adduced by commentators include Prov 26:11, %
and the uncleanness of dogs as scavengers in Jewish and Greek literature is a commonplace.\endnote{%
Neyrey 2021:221--222.} But this passage from Philumenos has apparently gone overlooked in the scholarly literature. %
The influence of his remedy for rabies, %
though worthy of ridicule in an early Jehudite-Christian social context, %
may place the composition/redaction of 2Peter in the early to mid-3rd century. %


